% ============================================================================
% Fichier cible : ARIEL_LOGMINER_MEMOIRE_FINAL/chapters/chapitre2_etat_art.tex
% Passage à remplacer : de « \section{Normalisation de journaux hétérogènes} »
%   (inclus) à « \section{Analyse séquentielle des journaux} » (exclu).
% Agent : Agent 3 — État de l'art et positionnement, chantier 2026-11-05.
% Sources : spécifications officielles ECS et OCSF ; résultats multiformat du corps.
% ============================================================================

\section{Normalisation de journaux hétérogènes}

\subsection{Du format source au contrat commun}

Le parsing et la normalisation répondent à deux questions distinctes. Le parsing extrait
la structure d'une source ; la normalisation projette les informations extraites dans un
vocabulaire commun afin que des règles, des recherches et des visualisations puissent
être réutilisées. Un contrat local peut remplir cette fonction dans un prototype, mais il
doit être comparé aux schémas déjà employés par les plateformes de sécurité.

L'\textit{Elastic Common Schema} (ECS) définit des noms de champs, leurs types et des
conventions d'extension pour des événements issus de sources variées. Son objectif est de
permettre l'analyse, la visualisation et la corrélation de données normalisées dans
l'écosystème Elastic \cite{elastic2026ecs}. L'\textit{Open Cybersecurity Schema
Framework} (OCSF) propose pour sa part un schéma de sécurité indépendant du fournisseur,
structuré autour d'une taxonomie, de classes d'événements, d'objets et d'un dictionnaire
d'attributs ; il est indépendant du format de stockage et du processus d'ingestion
\cite{ocsf2026schema}. ECS est ainsi étroitement articulé à une pile d'analyse donnée,
tandis qu'OCSF vise d'abord l'interopérabilité sémantique entre producteurs et
consommateurs.

\begin{table}[H]
\centering
\caption{Schémas de normalisation et contrat local}
\label{tab:schemas-normalisation}
\small
\begin{tabularx}{\textwidth}{p{2.3cm}p{3.2cm}p{3.5cm}X}
\toprule
\textbf{Schéma} & \textbf{Portée} & \textbf{Mécanisme principal} &
\textbf{Portée pour le prototype} \\
\midrule
ECS \cite{elastic2026ecs} & Événements, journaux et métriques analysés dans la pile Elastic. &
Champs typés, niveaux de champ et conventions pour les extensions. &
Référence concrète pour nommer les attributs usuels et réutiliser du contenu analytique. \\
OCSF \cite{ocsf2026schema} & Événements de cybersécurité indépendants du fournisseur et du stockage. &
Taxonomie, classes, objets, attributs, profils et extensions. &
Référence pour distinguer le noyau commun d'une taxonomie complète et versionnée. \\
Contrat local & Enveloppe interne du prototype. &
Petit noyau de champs et conservation d'attributs spécifiques ou de l'événement brut. &
Facilite les échanges internes, sans établir une conformité ECS ou OCSF. \\
\bottomrule
\end{tabularx}
\end{table}

\subsection{Extension, conservation et preuve de conformité}

Une normalisation trop agressive supprime des attributs propres à une source ; une
normalisation trop permissive produit un champ commun nominalement présent mais vide ou
sans sémantique stable. ECS comme OCSF prévoient des mécanismes d'extension, ce qui
confirme qu'un schéma commun ne doit pas remplacer l'événement source. Conserver le brut
ou une zone d'extension reste nécessaire pour l'audit et pour les traitements
spécialisés.

La conformité doit donc être évaluée champ par champ : présence, type, sémantique,
provenance et taux de remplissage. La simple émission d'une enveloppe ne prouve ni la
qualité du mapping ni l'interopérabilité avec un standard. Dans le protocole multiformat
de ce mémoire, les flux tabulaires ne renseignent aucun des sept champs standard retenus.
Le résultat démontre la lecture et l'émission d'un format interne sur les sources testées,
mais pas une normalisation complète au sens d'ECS ou d'OCSF. Le positionnement scientifique
est par conséquent celui d'une \textbf{projection commune partielle et extensible}, dont
la couverture doit être mesurée, et non celui d'un nouveau standard de normalisation.

